\section{强化学习关键现象：长训练、熵塌缩与探索维持}

Light SFT 为 RL 提供可用起点，但并不保证长期学习。随着策略集中到少数高奖励轨迹，entropy、任务难度和样本多样性会共同决定训练是否继续产生新能力；本节把这三个现象作为动态监控对象，而不是在训练结束后只看最终 reward。

\subsection{为什么 ``训练更久'' 不是自动有效}
增加 optimizer step 只有在策略仍能探索且 verifier 仍提供区分信号时才有价值。本小节分析 reward 上升与 entropy 下降可能同时发生的情况，并说明平台期可能来自策略变窄，而非任务已经被真正解决。
讲者提出第一个观察：很多系统在 RL 中会很快进入性能平台期。根因之一是输出熵下降过快，模型只会重复高概率轨迹，难以继续发现新解。对长任务 Agent 来说，这意味着 ``能做一点，但做不深''。

\begin{importantbox}{熵塌缩的实质}
熵塌缩不是 ``模型变差''，而是 ``模型变窄''：它在少数轨迹上越来越确定，但失去探索不同正确路径的能力。对于复杂任务，这会直接限制上限。
\end{importantbox}

\subsection{第二观察：任务难度需要动态调度}
如果训练样本长期停留在简单任务，模型很快 ``学完''，后续更新收益极低；但如果一开始就喂过难任务，也会因为正反馈稀疏而学不动。因此，课程隐含了一个 curriculum 观点：难度应随能力上升动态调整。

\begin{knowledgebox}{难度调度的实践框架}
\begin{itemize}
    \item 起步阶段：高可验证、低复杂度任务保证有效学习信号。
    \item 中期阶段：提高组合复杂度，增加多工具协作场景。
    \item 后期阶段：引入长轨迹、弱监督、开放式问题，重点保探索。
\end{itemize}
\end{knowledgebox}

\subsection{第三观察：多样高质量响应是核心资源}
课程把 ``diverse high-quality responses'' 作为第三支柱。这里的多样不是随机噪声，而是多条可行策略路径。只有当正负样本都具有结构多样性，RL 才能学到 ``什么时候该怎么做''，而不是死记模板。

\begin{warningbox}{伪多样性的陷阱}
如果多样性只体现在表面措辞，而动作序列和决策逻辑仍然单一，训练效果会被高估。真正有效的多样性应体现在工具选择、步骤顺序、失败恢复策略等决策层。
\end{warningbox}

\subsection{本章小结}
长训练要有效，必须同时处理三件事：防熵塌缩、做难度调度、保证策略级多样性。否则训练越久，收益越小。

